\section{选点模型}\label{sec:model}

\subsection{源的不确定集}

设第一检测点 $S_1$（取为原点）、示向度读数 $\theta_1$（取为 $0^\circ$）、测向误差上界
$\varepsilon=1^\circ$。一次测向只排除方向不合的位置，源到 $S_1$ 的距离 $d$ 未定，但有两个
硬边界：源必在 $1800$ m 圆域内；$S_1$ 处测出了示向度，故 $d>5$ m 且 $d\le1500$ m。于是源
的不确定集为

\begin{equation}
  C=\Bigl\{\,G:\;d=|S_1G|\in[5,1500]\ \text{m},\;
  \bigl|\angle(S_1\!\to\! G)-\theta_1\bigr|\le\varepsilon\,\Bigr\}\cap D(1800),
  \label{eq:C}
\end{equation}

即一个以 $S_1$ 为顶点、深 $1500$ m、张角 $2^\circ$ 的窄楔形。用 $W(p,\psi)$ 表示顶点 $p$、
中心方向 $\psi$、半张角 $\varepsilon$ 的楔形，$D(R)$ 表示半径 $R$ 的圆域，则
$C=W(S_1,\theta_1)\cap D(1500)\cap D(1800)$。

$C$ 是凸的（楔形与圆域皆凸），其极点只有 $4$ 个：$d\in\{5,1500\}$ 与
$\angle=\theta_1\pm\varepsilon$ 的四个组合。灵敏度分析表明 $d=5$ m 一侧对最优解无影响，
起决定作用的是两个远端极点 $d=1500$ m、$\angle=\theta_1\pm1^\circ$。

在算法中 $C$ 被离散为 $31\times9=279$ 个采样点（$d$ 步长 $50$ m、角向步长 $0.25^\circ$），
四个极点必须包含在采样内——后文将看到最坏情形恰好出现在极点与误差端点处。

\subsection{可测性保证与可行域}

设源的接收半径 $\rho_{\mathrm{rec}}\in[1000,1500]$ m 未知。要\textbf{保证}第二次测向成功，
必须对 $C$ 内每一个可能位置 $G$ 都有 $|S_2G|\le\rho_{\mathrm{rec}}$。由于
$\rho_{\mathrm{rec}}$ 的下界是 $1000$ m，最保守的保证型约束为

\begin{equation}
  \max_{G\in C}\,|S_2G|\;\le\;1000\ \text{m}.
  \label{eq:feasible}
\end{equation}

距离函数对 $S_2$ 是凸的，而 $C$ 是紧凸多边形，因此 $\max_{G\in C}|S_2G|$ 必在 $C$ 的极点
取得。记 $C$ 的四个极点为 $c_1,\dots,c_4$，约束~\eqref{eq:feasible} 等价于

\begin{equation}
  S_2\in F:=\bigcap_{i=1}^{4}D(c_i,\,1000),
  \label{eq:F}
\end{equation}

即四个半径 $1000$ m 的圆盘之交，是一个凸透镜形区域。按题面数据计算，
$|F|=0.4465$ km$^2$（占 $1800$ m 圆域的 $0.44\%$，占源不确定集楔形的 $0.13\%$），并且
$F\subset D(1800)$ 自动成立。$F$ 在极坐标下的径向区间可以由圆盘的射线距离解析给出，
后文构造候选弧带时直接使用该解析式。

文献中常用的处理是把"能否听到"写成软代价（听不到就再跑一趟），本文采用最保守的
"保证"口径；若放宽到接收半径上界 $1500$ m，可行域扩到 $2.737$ km$^2$、最优最坏直径
降到 $79.99$ m，该对照见 §\ref{sec:sens}。

\subsection{定位区域与最坏定位直径}

给定 $S_2$，设源在 $G$、源处第二次测向的名义方向为 $\theta_2=\angle(S_2\!\to\!G)$，读数
再叠加误差 $\delta_2\in[-\varepsilon,\varepsilon]$。与问题一一致，定位区域取两条示向度
误差边界射线围成的凸多边形（被 $1800$ m 圆域裁剪）：

\begin{equation}
  R(S_1,S_2;G,\delta_2)=W(S_1,\theta_1)\cap W\bigl(S_2,\theta_2+\delta_2\bigr)\cap D(1800).
  \label{eq:region}
\end{equation}

定位精度用区域\textbf{直径}度量：

\begin{equation}
  J(S_2)=\max_{G\in C}\ \max_{\delta_2\in[-\varepsilon,\varepsilon]}
  \ \mathrm{diam}\bigl(R(S_1,S_2;G,\delta_2)\bigr),
  \label{eq:J}
\end{equation}

其中 $\mathrm{diam}(\cdot)$ 为集合的最大点距。式~\eqref{eq:J} 的两层 $\max$ 分别对应
"源在楔形内何处"与"第二次测向偏多少"，两者都必须取最坏——这正是"保证"与"平均"的分水岭。
由于 $R$ 是凸多边形，其直径必在顶点对之间取得，故 $\mathrm{diam}$ 只需枚举顶点对（一般为
$4$ 个顶点，共 $6$ 对）即可精确求得，无需描点搜索。当两条射线近共线、区域退化为极窄条带时，
程序返回哨兵值（数值上等价于"无法定位"），保证统计与出图均为有限量。

\subsection{选点模型与候选区域}

选点模型即

\begin{equation}
  S_2^*=\operatorname*{arg\,min}_{S_2\in F}\ J(S_2),
  \qquad J^*=J(S_2^*).
  \label{eq:opt}
\end{equation}

为回答"周边区域适合程度"，定义适合度

\begin{equation}
  \mathcal{F}(S_2)=\frac{J^*}{J(S_2)}\in(0,1],\qquad
  \mathcal{F}(S_2^*)=1,
  \label{eq:Fitness}
\end{equation}

即 $J$ 越小越接近 $1$。候选区域取适合度阈值集

\begin{equation}
  \mathcal{C}(\eta)=\Bigl\{\,S_2\in F:\ \mathcal{F}(S_2)\ge\frac{1}{1+\eta}\,\Bigr\}
  =\Bigl\{\,S_2\in F:\ J(S_2)\le(1+\eta)\,J^*\,\Bigr\},
  \label{eq:cand}
\end{equation}

本文取 $\eta=10\%$，即"最坏直径不超过最优值 $1.1$ 倍"。$\eta$ 越小弧带越窄（§\ref{sec:sens}
给出 $\eta=5\%$ 与 $2\%$ 的结果）。

\subsection{一阶解析式：最优解为何"远且偏"}

把式~\eqref{eq:region} 的 $4$ 个顶点对源的真实方向做一阶（小角）展开，得到以 $2A$、$2B$ 为
邻边的平行四边形，其中

\begin{equation}
  A=\frac{r_1\varepsilon}{\sin\gamma},\qquad
  B=\frac{r_2\varepsilon}{\sin\gamma},\qquad
  \gamma=\text{源处两条示向度的交会角},
  \label{eq:AB}
\end{equation}

其对角线即直径的一阶近似：

\begin{equation}
  \mathrm{diam}\ \approx\ \frac{2\varepsilon}{\sin\gamma}\,
  \sqrt{\,r_1^2+r_2^2+2\,r_1r_2\,|\cos\gamma|\,}.
  \label{eq:closed}
\end{equation}

式~\eqref{eq:closed} 一眼说明三件事：\textbf{交会角越接近 $90^\circ$ 越好}
（$1/\sin\gamma$ 在 $\gamma=90^\circ$ 最小）；\textbf{两个检测点离源越近越好}（$r_1,r_2$ 越小
越好）；\textbf{三点近共线时直径发散}（$\gamma\to0^\circ$ 或 $180^\circ$），这正是必须避开的
退化构型。图~\ref{fig:criteria}(a) 给出该式的曲线（以精确构造的点作对照），可见 $1/\sin\gamma$
是主导项，与文献的几何稀释结论完全一致（§\ref{sec:lit}）。

式~\eqref{eq:closed} 只是趋势工具，\textbf{选点一律走精确构造}：400 例抽样中，解析构造与
shapely 精确求交的最大相对偏差为 $1.2\times10^{-14}$（同一套几何）；而式~\eqref{eq:closed}
与精确值在 $\gamma\in[30^\circ,120^\circ]$ 上的中位偏差为 $0.078\%$、$98\%$ 的样本小于
$1\%$，但在长条构型（$\gamma$ 很小、或 $r_2/d$ 很大）可达 $36\%$ 甚至 $65\%$，在本模型
最优点上给出 $121.49$ m、比精确值 $134.08$ m 低 $9.4\%$（二阶项与误差同向叠加的残余）。
因此它是解释性结论，不是计算路径。
